\section{NLO cut masters and endpoint reconstruction}
\label{sec:nlo-example}

The NLO real-emission sector is simple enough to solve fully while retaining
the same cuts, branches, and normalizations needed at higher orders.  It also
shows why the number of masters, the number of denominator classes, and the
number of independent boundary periods are different quantities.

\subsection{Two-particle cut and six-master basis}

Let $k_a,k_b,k_c$ be future-directed massless parton momenta and
$q=k_a+k_b-k_c$.  The invariants and chamber are those of
Eqs.~\eqref{eq:partonic-invariants}--\eqref{eq:physical-chamber}.  With one
independent emitted momentum $k$, define
\begin{equation}
 \dd\Phi_2(q;k,q-k)
 =(2\pi)^{2-D}\dd^Dk\,
 \delta_+(k^2)\delta_+((q-k)^2).
 \label{eq:nlo-two-body-measure}
\end{equation}
Its integral is
\begin{equation}
 \Phi_2(Q^2,\eps)
 =\frac{(4\pi)^\eps}{8\pi}
  \frac{\Gamma(1-\eps)}{\Gamma(2-2\eps)}(Q^2)^{-\eps}.
 \label{eq:nlo-two-body-volume}
\end{equation}
It is convenient to remove
\begin{equation}
 \mathcal P_2(Q^2,\eps)
 =\frac{(4\pi)^\eps\Gamma(1-\eps)}
        {16\pi^2\Gamma(1-2\eps)}(Q^2)^{-\eps},
 \qquad
 \Phi_2=\frac{2\pi}{1-2\eps}\mathcal P_2.
 \label{eq:nlo-reduced-phase-space-factor}
\end{equation}
No hidden factor $\mu^{2\eps}$ is included in either expression.

For $\sigma_i=\pm1$ and external vectors $R_i$, introduce
\begin{align}
 \widehat B(\eps)
 &=\mathcal P_2^{-1}\int\dd\Phi_2,
 \label{eq:nlo-bubble-definition}\\
 \widehat T[\sigma_1,R_1;\sigma_2,R_2]
 &=\mathcal P_2^{-1}\int\dd\Phi_2\,
 \frac{1}{(\sigma_1k+R_1)^2(\sigma_2k+R_2)^2}.
 \label{eq:nlo-top-definition}
\end{align}
The signs $\sigma_i$ specify momentum flow, not an $i0$ prescription.
The ordinary denominators have no interior zero in the chamber
\eqref{eq:physical-chamber}; their two boundary values coincide.

The six reduced masters are $\widehat B$ and
\begin{equation}
\begin{aligned}
 \widehat T_{ac}&=\widehat T[-1,k_a;+1,k_c],
 &\widehat T_{ab}&=\widehat T[-1,k_a;-1,k_b],\\
 \widehat T_{abb}&=\widehat T[-1,k_a;-1,k_a-k_c],
 &\widehat T_{cb}&=\widehat T[+1,k_c;-1,k_b],\\
 \widehat T_{cbb}&=\widehat T[+1,k_c;-1,k_a-k_c].
\end{aligned}
\label{eq:nlo-five-top-masters}
\end{equation}
They form three powered-integral classes and three denominator classes.  The
cut-only bubble class is contained in each of two maximal nontrivial
denominator classes, giving
\begin{equation}
 N_{\rm master}^{\rm NLO}=6,
 \qquad N_{\rm pow}^{\rm NLO}=3,
 \qquad N_{\rm den}^{\rm NLO}=3,
 \qquad N_{\rm max}^{\rm NLO}=2.
 \label{eq:nlo-class-counts}
\end{equation}

\subsection{One generic top integral}

Write $L_j(k)=\sigma_jk+R_j$ and define
\begin{equation}
 \Lambda_j=R_j^2+\sigma_jq\cdot R_j,
 \qquad
 \Delta_j=(q\cdot R_j)^2-q^2R_j^2.
 \label{eq:nlo-lambda-delta}
\end{equation}
For every line in Eq.~\eqref{eq:nlo-five-top-masters}, either $R_j^2=0$ or
$(q+\sigma_jR_j)^2=0$.  Hence $\Lambda_j^2=\Delta_j$, and the angular
dependence is described by one cross ratio
\begin{equation}
 z=\frac{1+\chi_{12}}{2},
 \qquad
 \chi_{12}=\sigma_1\sigma_2
 \frac{(q\cdot R_1)(q\cdot R_2)-q^2R_1\cdot R_2}
      {\Lambda_1\Lambda_2}.
 \label{eq:nlo-cross-ratio}
\end{equation}
The physical chamber implies $0<z<1$.

After the two cuts fix the energy and radial magnitude of $k$, the bubble
and top reduce to
\begin{align}
 \widehat B(\eps)&=\frac{2\pi}{1-2\eps},
 \label{eq:nlo-bubble-result}\\
 \widehat T(z,\eps;\Lambda_1,\Lambda_2)
 &=\frac{\pi}{\Lambda_1\Lambda_2}
 \int_0^\infty\dd y\,
 \frac{(1+y)^\eps}{1+(1-z)y}
 \notag\\
 &=-\frac{\pi}{\eps\Lambda_1\Lambda_2}
 {}_2F_1(1,1;1-\eps;z).
 \label{eq:nlo-generic-top}
\end{align}
The integral representation converges for $\operatorname{Re}\eps<0$.  Since
$1+(1-z)y>0$ for $y>0$ and $0<z<1$, it fixes a real branch before
meromorphic continuation in $\eps$.

The five physical top masters are obtained from one function by the
substitutions
\begin{table}[htbp]
\centering
\small
\begin{tabular}{@{}ccl@{}}
\toprule
Integral & $z_\alpha$ & $(\Lambda_{\alpha,1},\Lambda_{\alpha,2})$\\
\midrule
$\widehat T_{ac}$
& $\displaystyle\frac{su}{(s+t)(t+u)}$
& $\displaystyle\left(-\frac{s+t}{2},-\frac{t+u}{2}\right)$\\[2mm]
$\widehat T_{ab}$
& $\displaystyle\frac{tu}{(s+t)(s+u)}$
& $\displaystyle\left(-\frac{s+t}{2},-\frac{s+u}{2}\right)$\\[2mm]
$\widehat T_{abb}$
& $\displaystyle\frac{sQ^2}{(s+t)(s+u)}$
& $\displaystyle\left(-\frac{s+t}{2},-\frac{s+u}{2}\right)$\\[2mm]
$\widehat T_{cb}$
& $\displaystyle\frac{st}{(s+u)(t+u)}$
& $\displaystyle\left(-\frac{t+u}{2},-\frac{s+u}{2}\right)$\\[2mm]
$\widehat T_{cbb}$
& $\displaystyle\frac{uQ^2}{(s+u)(t+u)}$
& $\displaystyle\left(-\frac{t+u}{2},-\frac{s+u}{2}\right)$\\
\bottomrule
\end{tabular}
\caption{The five NLO top masters as substitutions into the generic angular
integral in Eq.~\eqref{eq:nlo-generic-top}.}
\label{tab:nlo-top-maps}
\end{table}

\subsection{Differential equation and boundary value}

Normalize the bubble to unity and remove the propagator scales,
\begin{equation}
 b=\frac{1-2\eps}{2\pi}\widehat B=1,
 \qquad
 h(z,\eps)=\frac{\Lambda_1\Lambda_2}{\pi}\widehat T.
 \label{eq:nlo-normalized-bubble-top}
\end{equation}
Using $\tau=y/(1+y)$, one has
\begin{equation}
 h(z,\eps)=\int_0^1\frac{\dd\tau\,(1-\tau)^{-1-\eps}}
 {1-z\tau}.
 \label{eq:nlo-euler-integral}
\end{equation}
The integral of the total derivative
\begin{equation}
 \frac{\dd}{\dd\tau}
 \left[\frac{\tau(1-\tau)^{-\eps}}{1-z\tau}\right]
 \label{eq:nlo-angular-total-derivative}
\end{equation}
vanishes for $\operatorname{Re}\eps<0$.  It yields the triangular system
\begin{equation}
 \frac{\partial}{\partial z}
 \begin{pmatrix}b\\h\end{pmatrix}
 =
 \begin{pmatrix}
 0&0\\[1mm]
 \dfrac{1}{z(1-z)}&\dfrac{z+\eps}{z(1-z)}
 \end{pmatrix}
 \begin{pmatrix}b\\h\end{pmatrix}.
 \label{eq:nlo-triangular-system}
\end{equation}
The integrating factor $z^{-\eps}(1-z)^{1+\eps}$ gives
\begin{equation}
 h=z^\eps(1-z)^{-1-\eps}
 \left[C(\eps)+\int_0^z\dd\xi\,
 \xi^{-1-\eps}(1-\xi)^\eps\right].
 \label{eq:nlo-general-de-solution}
\end{equation}
The convergent cut integral has
$h(0,\eps)=-1/\eps$.  Its homogeneous term behaves as
$C(\eps)z^\eps$, which diverges as $z\to0^+$ for
$\operatorname{Re}\eps<0$.  Therefore
\begin{equation}
 C(\eps)=0,
 \qquad
 h(z,\eps)=-\frac1\eps{}_2F_1(1,1;1-\eps;z).
 \label{eq:nlo-physical-de-solution}
\end{equation}
The five top masters thus need one generic boundary period.  Including the
elementary bubble volume,
\begin{equation}
 N_{\rm analytic\ inputs}^{\rm NLO}=2,
 \qquad
 N_{\rm non\text{-}elementary\ boundary\ periods}^{\rm NLO}=1.
 \label{eq:nlo-boundary-count}
\end{equation}

\subsection{Connection formula and endpoint branches}

Let $\rho=1-z$.  The exact connection formula is
\begin{align}
 h(z,\eps)={}&-\frac{1}{1+\eps}
 {}_2F_1(1,1;2+\eps;\rho)
 \notag\\
 &-\frac{\Gamma(1-\eps)\Gamma(1+\eps)}{\eps}
 z^\eps\rho^{-1-\eps}.
 \label{eq:nlo-connection-formula}
\end{align}
Both $z$ and $\rho$ are positive in the physical interval, so the regular
and singular modes in Eq.~\eqref{eq:nlo-connection-formula} are separated on
the real branch without a power-expansion rule.

Under the map in Eq.~\eqref{eq:vw-inverse}, the five cross ratios become
\begin{equation}
\begin{aligned}
 z_{ac}&=\frac{1-x}{1-vx},
 &z_{ab}&=\frac{(1-v)(1-x)}{1-v+vx},
 &z_{abb}&=\frac{x}{1-v+vx},\\
 z_{cb}&=\frac{1-v}{(1-v+vx)(1-vx)},
 &z_{cbb}&=\frac{v^2x(1-x)}{(1-v+vx)(1-vx)}.
\end{aligned}
\label{eq:nlo-cross-ratios-endpoint}
\end{equation}
Thus $z_{ac},z_{ab},z_{cb}\to1$, whereas
$z_{abb},z_{cbb}\to0$.  The common phase-space factor contributes
$(Q^2)^{-\eps}=(sv)^{-\eps}x^{-\eps}$.  Consequently the individual
masters contain the modes
\begin{equation}
 x^{-\eps},
 \qquad x^{-1-2\eps}.
 \label{eq:nlo-individual-master-modes}
\end{equation}
Rational hard-scattering coefficients may shift the integer part of an
exponent but not its coefficient of $\eps$.  After the complete diagram sum,
the singular real coefficient has the form
\begin{equation}
 \mathcal R_{\rm sing}(v,x,\eps)
 =C_1(v,\eps)x^{-1-\eps}
 +C_2(v,\eps)x^{-1-2\eps}.
 \label{eq:nlo-two-endpoint-branches}
\end{equation}
Applying Eq.~\eqref{eq:endpoint-distribution-identity} gives
\begin{align}
 \mathcal R_{\rm sing}={}&
 -\delta(1-w)\left(\frac{C_1}{\eps}+\frac{C_2}{2\eps}\right)
 +(C_1+C_2)\mathcal D_0(w)
 \notag\\
 &-\eps(C_1+2C_2)\mathcal D_1(w)
 +\frac{\eps^2}{2}(C_1+4C_2)\mathcal D_2(w)+\cdots .
 \label{eq:nlo-two-branch-distribution}
\end{align}
The Laurent expansions of $C_1$ and $C_2$ are inserted only after this
step.  In particular, a $1/\eps$ pole in either coefficient produces a
$1/\eps^2$ contribution multiplying $\delta(1-w)$.

\subsection{Combination with the remaining NLO terms}

Reverse unitarity determines the real contribution only.  In the convention
used by the existing UU calculation, the singular NLO coefficient is assembled
as
\begin{equation}
 H_{\mathsf{UU}}^{(1)}
 =H_{\mathsf{UU},\rm real}^{(1)}
 +H_{\mathsf{UU},\rm virt+UV}^{(1)}
 -H_{\mathsf{UU},\rm coll}^{(1)},
 \label{eq:nlo-uu-assembly}
\end{equation}
where $H_{\rm coll}$ denotes the positive convolution kernel that is
subtracted in the chosen factorization convention.  The stored analytic
calculation gives exact zero for the differences between the resulting
$\mathcal D_1$, $\mathcal D_0$, and $\delta(1-w)$ coefficients and the
independent unpolarized result of Ref.~\cite{Aversa:1988vb}.  This comparison
fixes the relative signs and normalization of the real, virtual, and collinear
terms in the UU channel.  The full regular remainder and the corresponding
LL and TT fixed-order combinations are separate tasks; they are not inferred
from this distributional comparison.
